\documentclass[11pt,addpoints]{exam}

\usepackage[margin=1in]{geometry}
\usepackage{amsmath, amssymb}
\usepackage{enumitem}
\usepackage{xcolor}

% Comment out the next line to produce the student version without solutions.
\printanswers

\renewcommand{\solutiontitle}{\noindent\textbf{Solution.}\enspace}
\unframedsolutions
\SolutionEmphasis{\color{blue!60!black}}

\pagestyle{headandfoot}
\firstpageheader{MATH 1920: Calculus II}{Midterm Exam 2}{October 2026}
\runningheader{MATH 1920}{Midterm 2 \ifprintanswers(Answer Key)\fi}{Page \thepage\ of \numpages}
\firstpagefooter{}{}{}
\runningfooter{}{}{}
\runningheadrule

\begin{document}

\begin{center}
  {\Large\bfseries Calculus II: Midterm Exam 2}\\[4pt]
  \ifprintanswers{\large\bfseries Answer Key}\\[4pt]\fi
  Integration Techniques, Improper Integrals, Sequences and Series\\[2pt]
  Time allowed: 75 minutes
\end{center}

\vspace{0.5em}
\makebox[\textwidth]{Name:\enspace\hrulefill\quad Student ID:\enspace\hrulefill}

\vspace{1em}
\noindent\textbf{Instructions.}
\begin{itemize}[nosep]
  \item Answer all questions. Show your work: correct answers without justification receive little or no credit.
  \item No calculators, notes or electronic devices are permitted.
  \item Simplify answers where reasonable, but you do not need to evaluate expressions such as $\ln 3$ or $\sqrt{7}$ numerically.
  \item The exam has \numquestions\ questions on \numpages\ pages, worth a total of \numpoints\ points.
\end{itemize}

\begin{center}
  \gradetable[h][questions]
\end{center}

\newpage

\begin{questions}

% ---------------------------------------------------------------
\question[10] Evaluate $\displaystyle \int x\,e^{2x}\,dx$.

\begin{solution}
Integrate by parts with $u = x$ and $dv = e^{2x}\,dx$, so $du = dx$ and $v = \tfrac{1}{2}e^{2x}$:
\[
  \int x e^{2x}\,dx = \frac{x}{2}e^{2x} - \int \frac{1}{2}e^{2x}\,dx = \frac{x}{2}e^{2x} - \frac{1}{4}e^{2x} + C = \frac{e^{2x}}{4}(2x - 1) + C.
\]
\end{solution}

% ---------------------------------------------------------------
\question[10] Evaluate $\displaystyle \int \frac{dx}{x^2\sqrt{9 - x^2}}$.

\begin{solution}
Let $x = 3\sin\theta$ with $-\frac{\pi}{2} < \theta < \frac{\pi}{2}$, $\theta \neq 0$. Then $dx = 3\cos\theta\,d\theta$ and $\sqrt{9 - x^2} = 3\cos\theta$, so
\[
  \int \frac{3\cos\theta\,d\theta}{9\sin^2\theta \cdot 3\cos\theta}
  = \frac{1}{9}\int \csc^2\theta\,d\theta
  = -\frac{1}{9}\cot\theta + C.
\]
From the reference triangle, $\cot\theta = \dfrac{\sqrt{9 - x^2}}{x}$. Therefore
\[
  \int \frac{dx}{x^2\sqrt{9 - x^2}} = -\frac{\sqrt{9 - x^2}}{9x} + C.
\]
\end{solution}

% ---------------------------------------------------------------
\question[12] Evaluate $\displaystyle \int \frac{2x^2 + x + 1}{(x + 1)(x^2 + 1)}\,dx$.

\begin{solution}
Decompose into partial fractions:
\[
  \frac{2x^2 + x + 1}{(x + 1)(x^2 + 1)} = \frac{A}{x + 1} + \frac{Bx + C}{x^2 + 1}.
\]
Multiplying through gives $2x^2 + x + 1 = A(x^2 + 1) + (Bx + C)(x + 1)$. Setting $x = -1$ gives $2 = 2A$, so $A = 1$. Then
\[
  x^2 + x = (Bx + C)(x + 1) = Bx^2 + (B + C)x + C,
\]
so $C = 0$ and $B = 1$. Hence
\[
  \int \left( \frac{1}{x + 1} + \frac{x}{x^2 + 1} \right) dx = \ln|x + 1| + \frac{1}{2}\ln(x^2 + 1) + C.
\]
\end{solution}

% ---------------------------------------------------------------
\question Determine whether each improper integral converges or diverges. If it converges, find its value.
\begin{parts}
  \part[6] $\displaystyle \int_1^\infty \frac{\ln x}{x^2}\,dx$
  \begin{solution}
  By parts with $u = \ln x$, $dv = x^{-2}\,dx$, so $du = dx/x$ and $v = -1/x$:
  \[
    \int \frac{\ln x}{x^2}\,dx = -\frac{\ln x}{x} + \int \frac{dx}{x^2} = -\frac{\ln x}{x} - \frac{1}{x} + C.
  \]
  Then
  \[
    \int_1^\infty \frac{\ln x}{x^2}\,dx = \lim_{t \to \infty}\left[ -\frac{\ln x + 1}{x} \right]_1^t
    = \lim_{t \to \infty}\left( -\frac{\ln t + 1}{t} \right) + 1 = 0 + 1 = 1,
  \]
  where the limit is $0$ by L'H\^opital's rule. The integral \textbf{converges to 1}.
  \end{solution}

  \part[4] $\displaystyle \int_0^2 \frac{dx}{(x - 1)^2}$
  \begin{solution}
  The integrand is unbounded at $x = 1$, inside the interval, so we must split the integral there. On the left piece,
  \[
    \int_0^1 \frac{dx}{(x - 1)^2} = \lim_{t \to 1^-}\left[ -\frac{1}{x - 1} \right]_0^t = \lim_{t \to 1^-}\left( \frac{1}{1 - t} - 1 \right) = \infty.
  \]
  Since one piece diverges, the integral \textbf{diverges}. (Ignoring the discontinuity and computing $\bigl[-\frac{1}{x-1}\bigr]_0^2 = -2$ is a common error; the result cannot be negative since the integrand is positive.)
  \end{solution}
\end{parts}

% ---------------------------------------------------------------
\question Determine whether each series converges absolutely, converges conditionally, or diverges. Name the test you use and verify its hypotheses.
\begin{parts}
  \part[3] $\displaystyle \sum_{n=1}^\infty \frac{n^2}{3^n}$
  \begin{solution}
  Ratio test: $\displaystyle \frac{a_{n+1}}{a_n} = \frac{(n+1)^2}{3^{n+1}} \cdot \frac{3^n}{n^2} = \frac{1}{3}\left(\frac{n+1}{n}\right)^2 \to \frac{1}{3} < 1$. \textbf{Converges absolutely.}
  \end{solution}

  \part[3] $\displaystyle \sum_{n=2}^\infty \frac{1}{n\ln n}$
  \begin{solution}
  Integral test: $f(x) = \frac{1}{x\ln x}$ is positive, continuous and decreasing for $x \ge 2$, and
  \[
    \int_2^\infty \frac{dx}{x\ln x} = \lim_{t \to \infty} \bigl[\ln(\ln x)\bigr]_2^t = \infty.
  \]
  \textbf{Diverges.}
  \end{solution}

  \part[3] $\displaystyle \sum_{n=1}^\infty \frac{(-1)^n}{\sqrt{n}}$
  \begin{solution}
  Alternating series test: $b_n = 1/\sqrt{n}$ is positive, decreasing and tends to $0$, so the series converges. But $\sum 1/\sqrt{n}$ is a $p$-series with $p = \frac12 \le 1$, which diverges. \textbf{Converges conditionally.}
  \end{solution}

  \part[3] $\displaystyle \sum_{n=1}^\infty \frac{n + 1}{n^3 + 2}$
  \begin{solution}
  Limit comparison with $b_n = 1/n^2$:
  \[
    \lim_{n \to \infty} \frac{(n + 1)/(n^3 + 2)}{1/n^2} = \lim_{n \to \infty} \frac{n^3 + n^2}{n^3 + 2} = 1,
  \]
  which is finite and positive. Since $\sum 1/n^2$ converges ($p = 2 > 1$), \textbf{the series converges absolutely} (its terms are positive).
  \end{solution}
\end{parts}

% ---------------------------------------------------------------
\question[12] Find the radius and interval of convergence of
\[
  \sum_{n=1}^\infty \frac{(x - 2)^n}{n\,3^n}.
\]

\begin{solution}
Ratio test on absolute values:
\[
  \lim_{n \to \infty} \left| \frac{(x - 2)^{n+1}}{(n+1)\,3^{n+1}} \cdot \frac{n\,3^n}{(x - 2)^n} \right| = \frac{|x - 2|}{3}\lim_{n \to \infty}\frac{n}{n + 1} = \frac{|x - 2|}{3}.
\]
The series converges when $|x - 2| < 3$, so the \textbf{radius is $R = 3$} and the open interval is $(-1, 5)$. Check the endpoints:
\begin{itemize}[nosep]
  \item $x = 5$: $\sum \frac{3^n}{n\,3^n} = \sum \frac{1}{n}$, the harmonic series, which diverges.
  \item $x = -1$: $\sum \frac{(-3)^n}{n\,3^n} = \sum \frac{(-1)^n}{n}$, the alternating harmonic series, which converges.
\end{itemize}
The \textbf{interval of convergence is $[-1, 5)$}.
\end{solution}

% ---------------------------------------------------------------
\question[10] Use a Maclaurin series to approximate $\displaystyle \int_0^1 e^{-x^2}\,dx$ with an error less than $0.001$. Justify the error bound.

\begin{solution}
From $e^u = \sum_{n=0}^\infty \frac{u^n}{n!}$ with $u = -x^2$,
\[
  e^{-x^2} = \sum_{n=0}^\infty \frac{(-1)^n x^{2n}}{n!},
  \qquad
  \int_0^1 e^{-x^2}\,dx = \sum_{n=0}^\infty \frac{(-1)^n}{n!\,(2n + 1)} = 1 - \frac{1}{3} + \frac{1}{10} - \frac{1}{42} + \frac{1}{216} - \frac{1}{1320} + \cdots
\]
This is an alternating series whose terms decrease in absolute value to $0$, so the error after stopping is less than the first omitted term. Since $\frac{1}{1320} \approx 0.00076 < 0.001$, five terms suffice:
\[
  \int_0^1 e^{-x^2}\,dx \approx 1 - \frac{1}{3} + \frac{1}{10} - \frac{1}{42} + \frac{1}{216} \approx 0.7475.
\]
(The exact value is $0.74682\ldots$, within the bound.)
\end{solution}

% ---------------------------------------------------------------
\question
\begin{parts}
  \part[6] The region bounded by $y = \sqrt{x}$ and $y = x^2$ is rotated about the $x$-axis. Find the volume of the resulting solid.
  \begin{solution}
  The curves meet where $\sqrt{x} = x^2$, at $x = 0$ and $x = 1$. On $[0, 1]$, $\sqrt{x} \ge x^2$. By the washer method,
  \[
    V = \pi\int_0^1 \left[ (\sqrt{x})^2 - (x^2)^2 \right] dx = \pi\int_0^1 (x - x^4)\,dx = \pi\left( \frac{1}{2} - \frac{1}{5} \right) = \frac{3\pi}{10}.
  \]
  \end{solution}

  \part[6] Find the arc length of the curve $y = \frac{2}{3}x^{3/2}$ for $0 \le x \le 3$.
  \begin{solution}
  Here $y' = x^{1/2}$, so $1 + (y')^2 = 1 + x$ and
  \[
    L = \int_0^3 \sqrt{1 + x}\,dx = \left[ \frac{2}{3}(1 + x)^{3/2} \right]_0^3 = \frac{2}{3}(8 - 1) = \frac{14}{3}.
  \]
  \end{solution}
\end{parts}

% ---------------------------------------------------------------
\question
\begin{parts}
  \part[6] Find the area enclosed by the cardioid $r = 1 + \cos\theta$.
  \begin{solution}
  The curve is traced once for $0 \le \theta \le 2\pi$:
  \[
    A = \frac{1}{2}\int_0^{2\pi} (1 + \cos\theta)^2\,d\theta = \frac{1}{2}\int_0^{2\pi} \left( 1 + 2\cos\theta + \frac{1 + \cos 2\theta}{2} \right) d\theta = \frac{1}{2}\left( 2\pi + 0 + \pi \right) = \frac{3\pi}{2}.
  \]
  \end{solution}

  \part[6] A curve is given parametrically by $x = t^2$, $y = t^3 - 3t$. Find all points where the tangent line is horizontal or vertical.
  \begin{solution}
  We have $\dfrac{dx}{dt} = 2t$ and $\dfrac{dy}{dt} = 3t^2 - 3 = 3(t - 1)(t + 1)$, so
  \[
    \frac{dy}{dx} = \frac{3(t^2 - 1)}{2t}.
  \]
  \begin{itemize}[nosep]
    \item \emph{Horizontal tangents} where $dy/dt = 0$ and $dx/dt \ne 0$: $t = \pm 1$, giving the points $(1, -2)$ at $t = 1$ and $(1, 2)$ at $t = -1$.
    \item \emph{Vertical tangent} where $dx/dt = 0$ and $dy/dt \ne 0$: $t = 0$, giving the point $(0, 0)$.
  \end{itemize}
  \end{solution}
\end{parts}

\end{questions}

\vfill
\begin{center}
  \textit{End of exam. Check that you have answered every part.}
\end{center}

\end{document}
